
\subsection{Model Architectures}
We evaluate XX distinct architectures designed to balance accuracy and computational efficiency.
The benchmark includes standard CNN-LSTM and ResNet variants, physics-informed PINN models, KAN heads, and continuous-time LTC-based temporal models.

\subsubsection{ResNet-based Architectures}
%SimpleResNet, SpatialResNet

\paragraph{Simple ResNet}
A baseline frame-based approach that predicts temperature from a single 5-channel image using a modified ResNet-18. It lacks explicit temporal memory (LSTM) but benefits from the strong spatial feature extraction of the ResNet architecture.

\subsubsection{CNN-LSTM Architectures}
%CNNLSTM, PretrainedCNNLSTM,

\paragraph{CNN-LSTM (Custom Lightweight)}

A compact architecture designed for extreme edge constraints. 
It consists of a 3-layer CNN feature extractor (16, 32, 64 filters) followed by a Global Average Pooling layer and a single LSTM layer (hidden size 64). 
This model has only $\sim$60k parameters.

\paragraph{Pretrained CNN-LSTM}

Leverages a ResNet-18 \cite{he2016deep} backbone pretrained on ImageNet. 
To accommodate the 5-channel input, we modify the first convolutional layer by extending its weights: the first 3 channels are initialized from the pretrained RGB weights, while the 2 flow channels are initialized using Kaiming initialization. 
The extracted features feed into an LSTM (hidden size 128).


\paragraph{Physics-Informed CNN-LSTM}
%PhysicsCNNLSTM, ConvectionBioheat

This architecture mirrors the Pretrained CNN-LSTM and adds a temporal regularization term based on Newton's law of cooling:
\begin{equation}
    \frac{dT}{dt} = -k(T - T_{env})
\end{equation}
where $T$ is the predicted temperature, $T_{env}$ is a reference temperature, and $k$ is a coefficient. The implemented residual uses frame-to-frame temperature differences. 
However, the experiments described in Section~\ref{sec:data_capture} record the heating phase under active MWA, not post-ablation cooling. 
This equation omits the microwave heat source and therefore does not describe the experimental heating process. 
We report this scalar model as using a phenomenological temporal regularizer, not as a mechanistic bioheat model.
The total loss function is defined as:
\begin{equation}
    \mathcal{L}_{total} = \mathcal{L}_{MSE} + \lambda \mathcal{L}_{physics}
\end{equation}
where $\mathcal{L}_{MSE}$ is the standard Mean Squared Error between predicted and ground truth temperatures, and $\mathcal{L}_{physics}$ penalizes residuals of the differential equation:
\begin{equation}
    \mathcal{L}_{physics} = \frac{1}{N} \sum_{i=1}^{N} \left| \frac{dT}{dt} + k(T - T_{env}) \right|^2
\end{equation}
This term regularizes the temporal predictions, but its physical validity for the heating-phase data is not assumed.

%\paragraph{Advanced Bioheat-Informed Model (Inverse PINN)}
The spatial model uses a reduced, Pennes-inspired residual with learnable coefficients $\alpha$ and $\beta$ and predicts a spatio-temporal temperature map $T(x,y,t)$.
In the implementation, the residual includes temporal change, spatial diffusion, a perfusion-like relaxation term, and a learnable constant source:
\begin{equation}
    \frac{\partial T}{\partial t} = \beta \nabla^2 T - \alpha(T - T_{arterial}) + Q_{met}
\end{equation}
This is a reduced surrogate rather than the full Pennes equation: it does not explicitly include the spatially distributed microwave heat source $Q_{\mathrm{MWA}}$. 
In addition, the agar phantoms have no blood perfusion, and the coefficients are not calibrated in physical units. 
Accordingly, $\alpha$, $\beta$, and $Q_{met}$ are interpreted as effective regularization parameters rather than measurements of tissue properties.
The loss minimizes the residual of this reduced equation; it does not establish that the predictions satisfy a mechanistic model of MWA heating.

It is important to note that the full application of this bioheat loss requires temporal sequence data to compute the time derivative $\frac{\partial T}{\partial t}$. 
When applied to single-frame architectures (e.g., Spatial ResNet), the loss function adapts by enforcing only the spatial diffusion term ($\beta \nabla^2 T$), effectively acting as a spatial smoothness regularizer without the temporal dynamics.
%XX

%\paragraph{Optical-Flow-Informed Surrogate}
To test whether apparent inter-frame image motion can provide an additional regularization signal, we add an advection-like term using dense optical flow $\vec{v} = (v_x, v_y)$ computed during preprocessing:
\begin{equation}
    \frac{\partial T}{\partial t} + \vec{v} \cdot \nabla T = \beta \nabla^2 T - \alpha(T - T_{arterial}) + Q_{met}
\end{equation}
The optical-flow vectors are downsampled to match the resolution of the predicted temperature map. 
Since optical flow from B-mode video is not a direct measurement of tissue velocity or fluid flow, and the agar phantoms have no perfusion, this term is an image-motion surrogate rather than a measurement of convection.

%\paragraph{Metabolic Heat Generation}
The loss includes a learnable constant source term $Q_{met}$:
\begin{equation}
    \frac{\partial T}{\partial t} + \vec{v} \cdot \nabla T = \beta \nabla^2 T - \alpha(T - T_{arterial}) + Q_{met}
\end{equation}
This constant is not separately identifiable as metabolic heating or microwave absorption.
The implemented residual does not explicitly include the spatial distribution of microwave energy deposition.

%\paragraph{Inverse Physics and Spatial Parameter Discovery}
Beyond temperature regression, our spatial model treats the effective coefficients $\alpha$ and $\beta$ as learnable spatial maps.
A parameter head predicts these $4\times4$ maps from the deep feature representation of the first frame.
Because the coefficients are uncalibrated and the phantom is avascular, they are model parameters for the reduced residual, not measurements of perfusion or thermal conductivity and not tissue-type labels.



\subsubsection{Kolmogorov--Arnold Heads (KAN-ResNet and Spatial KAN-Bioheat)}\label{sec:kan_models}
%KANResNet, SpatialKANBioheat

Motivated by the recent success of Kolmogorov--Arnold Networks (KANs) in scientific machine learning~\cite{liu2024kan,wang2024kinn,patra2024pikan} and by the observation that KANs benefit from a strong encoder when the input is high-dimensional~\cite{shen2025alpkan}, we introduce two architectures that retain our convolutional video encoder but replace the regression head with a KAN.

A KAN linear layer maps an input vector $\mathbf{x}\in\mathbb{R}^{d_{in}}$ to $\mathbf{y}\in\mathbb{R}^{d_{out}}$ via learnable univariate activations placed on each edge:
\begin{align}
    y_j &= \sum_{i=1}^{d_{in}} \bigl( w^{b}_{ji}\,\sigma(x_i) + w^{s}_{ji}\, \phi_{ji}(x_i) \bigr), \\
    \phi_{ji}(x_i) &= \sum_{k=1}^{G+p} c_{jik}\, B_{k,p}(x_i).
\end{align}
where $\sigma$ is a smooth base activation (SiLU), $B_{k,p}$ are B-splines of order $p$ over a uniform grid of $G$ knots, and $\{c_{jik}, w^{b}_{ji}, w^{s}_{ji}\}$ are the learnable parameters. 
We adopt the efficient implementation of~\cite{blealtan2024efficientkan}, which collapses the per-edge spline contraction into a single dense matrix multiply.

\paragraph{KAN-ResNet} 
A direct apples-to-apples ablation of the Simple ResNet baseline: the same 5-channel ResNet-18 backbone produces a 512-dimensional latent vector, which is fed to a two-layer KAN ($512 \to 64 \to 4$) that emits the four sensor temperatures. 
This isolates the effect of the head, holding the encoder fixed.

\paragraph{Spatial KAN-Bioheat}
A separable PIKAN field head in the style of SPIKANs~\cite{jacob2025spikan}. Three small one-dimensional KANs, $\Phi_x,\Phi_y,\Phi_t$, evaluate learnable basis functions on the spatial and temporal collocation coordinates of the $4\times4\times T$ output grid. 
Their outer product produces a coordinate feature tensor that is concatenated, voxel-wise, with the encoder latent and passed through a fusion KAN to produce the temperature field $T(x,y,t)$:
\begin{equation}
    T(x_i, y_j, t_k) \;=\; \mathrm{KAN}_{\text{fuse}}\!\Bigl(\, \Phi_t(t_k)\odot\Phi_y(y_j)\odot\Phi_x(x_i)\;\Vert\;\mathbf{z}_{\text{enc}} \,\Bigr).
\end{equation}
The field is trained with the same Advanced Bioheat Loss as our \emph{Convection-Bioheat} model, so that the comparison directly probes whether the spline-based functional form is a better surrogate than a tiny convolutional head for approximating the PDE residual on a coarse grid. 
Because the spline activations are smooth and differentiable to high order, the autograd-based residual $\partial_t T + \vec{v}\!\cdot\!\nabla T - \beta\nabla^2 T + \alpha(T-T_{art})$ is well-behaved without explicit numerical smoothing.

Both KAN variants enter the LOSO benchmark with no other change to the training pipeline, providing a controlled study of the encoder--KAN composition in a clinical ultrasound regression task.


\subsubsection{Convolutional Liquid Time-Constant Architecture}
\label{sec:conv_ltc}
Liquid Time-Constant networks provide a recurrent model with input-dependent
state dynamics~\cite{lechner2020neural}. The evaluated LTC model, ConvLTC,
consists of a two-layer convolutional
encoder, a convolutional Liquid Time-Constant (LTC) recurrent core, and a
convolutional decoder. The encoder maps each frame to feature maps while
preserving their spatial dimensions. The recurrent state likewise retains
the spatial layout.

At each time step, the recurrent cell concatenates the current encoded frame
with the previous hidden state. Convolutions generate an input-dependent
leak and drive signal, and an Euler update advances the hidden state. The
decoder maps each recurrent state to a temperature map, producing a
spatiotemporal sequence of predictions. In the current implementation, the
update uses a fixed unit time step; irregular sampling intervals are not
passed to the recurrent core.


%\subsubsection{Bayesian Architectures}

%BayesianResNet
%FullBayesianResNet
%BayesianCNNLSTM


\subsection{Uncertainty Quantification}
Reliability is paramount in clinical and industrial settings. To provide clinicians with confidence intervals alongside temperature estimates, we implement two distinct probabilistic approaches:

\subsubsection{Deep Ensembles}
We train $M=5$ independent instances of the SimpleResNet architecture. Each model is initialized with different random seeds and trained on the same dataset. During inference, the ensemble prediction is the mean of the individual outputs, and the variance represents the aleatoric uncertainty (data uncertainty). This method is computationally expensive at training time but offers a robust baseline for uncertainty estimation.

\subsubsection{Bayesian Neural Networks (BNN)}
We employ Variational Inference to approximate the posterior distribution of the network weights, $P(w|\mathcal{D})$, given the training data $\mathcal{D}$. Since the true posterior is intractable, we approximate it with a variational distribution $q_\theta(w)$ (typically a Gaussian) parameterized by $\theta$. We use the "Bayes by Backprop" algorithm \cite{blundell2015weight} to minimize the Evidence Lower Bound (ELBO), which is equivalent to minimizing the combined loss:
\begin{equation}
    \mathcal{L}_{BNN} = \mathcal{L}_{MSE} + \beta \cdot \text{KL}(q_\theta(w) || P(w))
\end{equation}
Note that because this objective is defined as a minimization loss, the KL term appears with a positive sign: we minimize the negative ELBO, which combines the data error with a regularizer that keeps the variational posterior close to the prior. Here, $\mathcal{L}_{MSE}$ is the data likelihood (reconstruction error), $\text{KL}$ is the Kullback-Leibler divergence between the variational posterior and the prior $P(w)$, and $\beta$ is a weighting factor balancing complexity and accuracy.

We investigate two BNN variants:
\begin{itemize}
    \item \textbf{Bayesian Head}: We use a standard ResNet-18 backbone for feature extraction and only replace the final fully connected layers with Bayesian layers. This approach is computationally efficient as the heavy convolutional backbone remains deterministic.
    \item \textbf{Full Bayesian}: We convert every learnable layer (Convolutional and Linear) in the ResNet-18 to a Bayesian layer. This captures epistemic uncertainty (model uncertainty) across the entire hierarchy of features but significantly increases the number of parameters (doubling them to store means and variances).
\end{itemize}
Implementation was performed using the \texttt{torchbnn} library. During inference, we perform multiple forward passes (Monte Carlo sampling) to estimate the predictive mean and variance.


We also evaluate a physics-informed Bayesian variant (B-PINN) proposing a hybrid architecture that synergizes the strengths of Bayesian Deep Learning and Physics-Informed Learning. 
This model employs a Bayesian ResNet backbone to capture epistemic uncertainty while simultaneously minimizing the Advanced Bioheat Loss. 
The total objective function becomes:
\begin{equation}
    \mathcal{L}_{total} = \mathcal{L}_{MSE} + \lambda_{phys} \mathcal{L}_{bioheat} + \lambda_{KL} \text{KL}(q_\theta(w) || P(w))
\end{equation}
This formulation ensures that the model's predictions are data-driven, physically consistent, and probabilistically calibrated. The physics loss acts as a regularizer that constrains the search space of the Bayesian posterior, potentially leading to better generalization in low-data regimes.
